\documentclass[11pt]{article}

\usepackage[a4paper,fancysections,titlepage]{src/polytechnique}
\usepackage[english]{babel}
\usepackage[T1]{fontenc}
\usepackage[hidelinks]{hyperref}
\usepackage{amsmath,amssymb}
\usepackage{booktabs}
\usepackage{array}
\usepackage{tikz}
\usetikzlibrary{arrows.meta, positioning, fit, backgrounds, shapes.geometric}
\usepackage{listings}
\usepackage{xcolor}
\usepackage{setspace}
\usepackage[patch=none]{microtype}

\setlength{\parskip}{0.7em}
\setstretch{1.05}

\lstset{
  basicstyle=\ttfamily\small,
  breaklines=true,
  frame=single,
  backgroundcolor=\color{bleu303pale},
  rulecolor=\color{bleu303},
  keywordstyle=\color{bleu315}\bfseries,
  commentstyle=\color{gray}\itshape,
  stringstyle=\color{rouge485},
  language=Python,
  numbers=left,
  numberstyle=\tiny\color{gray},
  xleftmargin=1.5em,
  framexleftmargin=1.5em,
}

\newcommand{\code}[1]{\mbox{\ttfamily\detokenize{#1}}}

\author{Nathan Duboisset}
\date{2026}
\title{PlaylistGenerator}
\subtitle{Semantic Music Retrieval from Natural Language and Visual Prompts}

\begin{document}

\maketitle

\tableofcontents

\newpage

% ─────────────────────────────────────────────────────────────────────────────
\section{Introduction}
% ─────────────────────────────────────────────────────────────────────────────

Music streaming platforms host hundreds of millions of tracks, yet their recommendation
engines still rely largely on collaborative filtering and rigid keyword tags.
\textbf{PlaylistGenerator} addresses this by building a semantic retrieval system that
maps audio and free-text queries into a shared embedding space using CLAP~\cite{elizalde2023clap}.

The system is built around four components:
\begin{enumerate}
    \item \textbf{Audio indexing} --- pre-compute CLAP audio embeddings for every track
          in FMA Small and store them in a compressed \texttt{.npz} index.
    \item \textbf{Text retrieval} --- encode the user prompt with the CLAP text encoder
          and run a cosine $k$-NN search against the index.
    \item \textbf{Two-phase fine-tuning} --- (i) adapt the text encoder with LoRA on
          MusicCaps; (ii) further fine-tune the last stage of the audio encoder
          end-to-end on raw waveforms to tighten genre clusters in audio embedding space.
    \item \textbf{Image-to-playlist bridge} --- a small MLP projects CLIP image features
          into the CLAP text space, enabling image-conditioned queries.
\end{enumerate}

% ─────────────────────────────────────────────────────────────────────────────
\section{Background}
% ─────────────────────────────────────────────────────────────────────────────

\subsection{CLAP}

CLAP~\cite{elizalde2023clap} is a dual-encoder model: the audio branch is an
\textbf{HTSAT-base} hierarchical Swin transformer (4 stages); the text branch is
\textbf{RoBERTa}. Both encode into a shared 512-dimensional space trained with the
symmetric InfoNCE loss:
\begin{equation}
    \mathcal{L} = \frac{1}{2}\Bigl(
        \mathcal{L}_{a\to t} + \mathcal{L}_{t\to a}
    \Bigr), \quad
    \mathcal{L}_{a\to t} = -\frac{1}{N}\sum_{i=1}^{N}
        \log\frac{e^{s\,\mathbf{a}_i^\top\mathbf{t}_i}}
                  {\sum_{j} e^{s\,\mathbf{a}_i^\top\mathbf{t}_j}}
    \label{eq:infonce}
\end{equation}
where $s = \exp(\texttt{logit\_scale})$ is learned and $N$ is the batch size.
The checkpoint used is \texttt{music\_audioset\_epoch\_15\_esc\_90.14.pt} (90.14\%
on ESC-50).

\subsection{LoRA and CLIP}

LoRA~\cite{hu2022lora} is used as-is to adapt the text encoder (see Section~\ref{sec:phase1}).
CLIP~\cite{radford2021clip} ViT-L/14 produces 1024-dim image embeddings that feed the
bridge MLP (Section~\ref{sec:bridge}).

% ─────────────────────────────────────────────────────────────────────────────
\section{Datasets}
% ─────────────────────────────────────────────────────────────────────────────

\textbf{FMA Small}~\cite{defferrard2017fma}: 8,000 tracks across 8 genres (Hip-Hop,
Pop, Folk, Experimental, Rock, International, Electronic, Instrumental), 30-second MP3s
at 44.1 kHz. Six corrupt files are silently skipped, giving \textbf{7,994 indexed tracks}.

\textbf{MusicCaps}~\cite{agostinelli2023musiclm}: 5,521 expert-captioned 10-second
YouTube clips. After downloading, 5,298 are available. Pairing with pre-computed audio
embeddings yields \textbf{4,357 samples} (85/15 split, seed 42: 3,704 train / 653 test).
Phase~2 uses raw WAV files directly rather than cached embeddings.

\textbf{MS COCO Captions}: 20,000 image--caption pairs for bridge MLP training. Each
image is encoded by frozen CLIP ViT-L/14; each caption by frozen CLAP text encoder.

% ─────────────────────────────────────────────────────────────────────────────
\section{Methods}
% ─────────────────────────────────────────────────────────────────────────────

\subsection{Overview}

Figure~\ref{fig:arch} shows the three inference paths. All paths converge on a cosine
$k$-NN search against a \texttt{.npz} index of 7,994 L2-normalised audio embeddings.

\begin{figure}[h]
\centering
\begin{tikzpicture}[
    font=\small,
    box/.style={draw, rounded corners=3pt, minimum width=2.6cm, minimum height=0.75cm,
                align=center, fill=bleu303pale, draw=bleu303},
    model/.style={draw, rounded corners=3pt, minimum width=2.6cm, minimum height=0.75cm,
                  align=center, fill=rouge485!15, draw=rouge485},
    store/.style={draw, cylinder, shape border rotate=90,
                  minimum width=1.8cm, minimum height=0.9cm,
                  align=center, fill=bleu315!15, draw=bleu315},
    arr/.style={-{Stealth[length=4pt]}, thick},
]
\node[box]  (textprompt)  at (0, 3)    {Text Prompt};
\node[model](claptext)    at (3, 3)    {CLAP\\Text Enc.\\(RoBERTa+LoRA)};
\node[box]  (imageprompt) at (0, 1.5)  {Image Prompt};
\node[model](clip)        at (3, 1.5)  {CLIP\\ViT-L/14};
\node[model](bridge)      at (6, 1.5)  {Bridge MLP\\$1024\!\to\!512$};
\node[box]  (audiocorpus) at (0, 0)    {FMA Small};
\node[model](clapaudio)   at (3, 0)    {CLAP Audio\\Enc. (HTSAT)};
\node[store](index)       at (6.2, 0)  {\texttt{.npz}\\$7994\times512$};
\node[model](knn)         at (9, 2.25) {Cosine $k$-NN};
\node[box]  (playlist)    at (12,2.25) {Playlist};
\draw[arr] (textprompt)  -- (claptext);
\draw[arr] (claptext)    -- node[above]{\small $e_t$} (9,3) -- (knn);
\draw[arr] (imageprompt) -- (clip);
\draw[arr] (clip) -- (bridge);
\draw[arr] (bridge) -- node[right]{\small $e_v$} (9,1.5) -- (knn);
\draw[arr] (audiocorpus) -- (clapaudio);
\draw[arr] (clapaudio)   -- (index);
\draw[arr] (index) -- node[above]{\small offline} (9,0) -- (knn);
\draw[arr] (knn) -- (playlist);
\end{tikzpicture}
\caption{System overview. The audio index is built offline; retrieval is a single
matrix--vector product ($\mathcal{O}(N)$).}
\label{fig:arch}
\end{figure}

\subsection{Phase 1 — LoRA Text Encoder Fine-tuning}
\label{sec:phase1}

The text encoder (RoBERTa) is adapted to music-specific captions via LoRA
(\texttt{r=16}, $\alpha=32$, dropout=0.1, targets: \texttt{query}, \texttt{value}),
adding \textbf{1,246,209 trainable parameters} (0.62\% of 200M). The audio encoder
is \textbf{fully frozen}; audio embeddings are pre-computed once and cached.

Each training step minimises the symmetric InfoNCE loss (Eq.~\ref{eq:infonce}) between
a batch of cached audio embeddings and freshly encoded text embeddings. Hyperparameters:
AdamW, $lr=10^{-4}$, weight decay $10^{-4}$, cosine schedule, batch 16, 3 epochs,
gradient clipping at 1.0.

\subsection{Phase 2 — Audio Encoder Fine-tuning}
\label{sec:phase2}

After Phase~1 the LoRA weights are frozen and the \textbf{last Swin stage}
(\texttt{audio\_branch.layers[3]}) together with the final norm,
\texttt{audio\_transform}, and \texttt{audio\_projection} are unfrozen — roughly
the top 15\% of the audio encoder by depth. This targeted strategy preserves
low-level spectral features while specialising high-level representations to music.

Because the audio encoder is now updated, pre-computed embeddings are stale.
Training instead loads raw 10-second WAV files, resampled on-the-fly to 48\,kHz
(CLAP's native rate), and passes them through the live audio encoder. The text
encoder runs in eval mode (no LoRA dropout) with no gradient. The loss is the same
InfoNCE objective; the LoRA text embeddings serve as fixed targets.

Hyperparameters: AdamW, $lr=5\times10^{-6}$ (conservative to avoid forgetting),
weight decay $10^{-4}$, cosine schedule, batch 8, 2 epochs, gradient clipping 1.0.

\subsection{Evaluation Metrics}

\paragraph{Text--audio alignment.} Mean cosine similarity between the CLAP text
embedding of each genre name and the mean audio embedding of its tracks
(computed on FMA Small). Measures how well text labels match the audio cluster centroid.

\paragraph{Intra-genre cohesion.} Mean \emph{pairwise} cosine similarity among all
audio embeddings within a genre:
\[
    C_g = \frac{1}{n_g(n_g-1)}\sum_{i \neq j} \mathbf{a}_i^\top \mathbf{a}_j
\]
A higher $C_g$ means tracks of the same genre are closer together in embedding space —
a direct measure of how well the audio encoder clusters by genre.
This metric is independent of the text encoder in Phase~1 but \emph{changes} in
Phase~2 since the audio encoder is updated.

\subsection{\texorpdfstring{CLIP$\to$CLAP}{CLIP-to-CLAP} Bridge MLP}
\label{sec:bridge}

A two-hidden-layer MLP maps 1024-dim CLIP image features to 512-dim CLAP embeddings
($1024\to1024\to1024\to512$, ReLU + LayerNorm + Dropout$_{0.2}$). Trained with
one-directional InfoNCE ($\tau=0.07$) on 20,000 MS COCO pairs; AdamW
$lr=2\times10^{-4}$, batch 64, 5 epochs.

% ─────────────────────────────────────────────────────────────────────────────
\section{Experiments and Results}
% ─────────────────────────────────────────────────────────────────────────────

\subsection{Phase 1 — LoRA Text Encoder}

\begin{table}[h]
\centering
\caption{InfoNCE loss per epoch — Phase~1 (MusicCaps, 3704/653 samples).}
\label{tab:training}
\begin{tabular}{ccc}
\toprule
\textbf{Epoch} & \textbf{Train} & \textbf{Val} \\
\midrule
1 & 0.6165 & 0.4093 \\
2 & 0.3959 & 0.3659 \\
3 & 0.3277 & 0.3428 \\
\bottomrule
\end{tabular}
\end{table}

\begin{table}[h]
\centering
\caption{Average cosine similarity on the MusicCaps test set (653 samples).}
\label{tab:cosine}
\begin{tabular}{lcc}
\toprule
\textbf{Model} & \textbf{Avg.\ cosine sim.} & \textbf{Gain} \\
\midrule
Baseline CLAP         & 0.3135 & --- \\
CLAP + LoRA (Phase~1) & 0.3710 & $+18.3\%$ \\
\bottomrule
\end{tabular}
\end{table}

LoRA improves text--audio alignment by \textbf{+0.0575 absolute} with only 0.62\%
of parameters updated.

\subsection{Per-genre Text--Audio Alignment}

Table~\ref{tab:genre_align} shows per-genre alignment before and after Phase~1.
Genres with low baseline alignment (Electronic, Experimental) see the largest gains;
genres already well-aligned (Folk, Hip-Hop) regress slightly, likely because
MusicCaps skews toward Electronic/Experimental styles.

\begin{table}[h]
\centering
\caption{Per-genre text--audio cosine similarity on FMA Small.}
\label{tab:genre_align}
\begin{tabular}{lccc}
\toprule
\textbf{Genre} & \textbf{Baseline} & \textbf{Phase 1} & \textbf{$\Delta$} \\
\midrule
Electronic    & 0.1085 & 0.2041 & $+0.0956$ \\
Experimental  & 0.2349 & 0.2856 & $+0.0507$ \\
Instrumental  & 0.2466 & 0.2706 & $+0.0240$ \\
Rock          & 0.3540 & 0.3271 & $-0.0269$ \\
Hip-Hop       & 0.4212 & 0.3471 & $-0.0741$ \\
Folk          & 0.4135 & 0.3133 & $-0.1001$ \\
Pop           & 0.2537 & 0.1383 & $-0.1155$ \\
International & 0.3483 & 0.1850 & $-0.1632$ \\
\bottomrule
\end{tabular}
\end{table}

\subsection{Intra-genre Cohesion and Phase 2}

Table~\ref{tab:cohesion} reports the intra-genre cohesion $C_g$ before Phase~2
(frozen audio encoder) and after Phase~2 (fine-tuned last Swin stage). Because
Phase~1 does not touch the audio encoder, the ``Baseline'' and ``Phase~1'' columns
are identical; only Phase~2 changes $C_g$.

\begin{table}[h]
\centering
\caption{Intra-genre audio cohesion $C_g$ on FMA Small.}
\label{tab:cohesion}
\begin{tabular}{lcc}
\toprule
\textbf{Genre} & \textbf{Phase 1 (frozen audio)} & \textbf{Phase 2 ($\Delta$)} \\
\midrule
Hip-Hop       & \emph{run notebook} & --- \\
Folk          & \emph{run notebook} & --- \\
Rock          & \emph{run notebook} & --- \\
Electronic    & \emph{run notebook} & --- \\
Experimental  & \emph{run notebook} & --- \\
Instrumental  & \emph{run notebook} & --- \\
Pop           & \emph{run notebook} & --- \\
International & \emph{run notebook} & --- \\
\bottomrule
\end{tabular}
\end{table}

An increase in $C_g$ after Phase~2 indicates that the fine-tuned audio encoder
maps same-genre tracks closer together, making $k$-NN retrieval more genre-coherent.

\subsection{Qualitative Retrieval}

Table~\ref{tab:qualitative} shows top-5 MusicCaps results for three prompts with
the Phase~1 model. Retrieved captions are semantically consistent with the query
in all three cases.

\begin{table}[h]
\centering
\caption{Top-5 retrieval from MusicCaps (Phase~1 model).}
\label{tab:qualitative}
\setlength{\tabcolsep}{4pt}
\begin{tabular}{clc}
\toprule
\textbf{Rank} & \textbf{Caption (truncated)} & \textbf{Score} \\
\midrule
\multicolumn{3}{l}{\textit{``slow chill hip hop with a mellow bassline''}} \\
\midrule
1 & Hip hop song with consistent low-quality bassline   & 0.405 \\
2 & Dark moody piano melody with hip hop beat           & 0.392 \\
3 & Instrumental movie soundtrack theme                 & 0.390 \\
\midrule
\multicolumn{3}{l}{\textit{``upbeat reggae with rhythmic percussion''}} \\
\midrule
1 & Acoustic drum playing a reggae groove with rim shot & 0.421 \\
2 & Female singer over a cool melody                    & 0.413 \\
3 & Reggae song with flat male vocal rapping            & 0.384 \\
\midrule
\multicolumn{3}{l}{\textit{``sad orchestral strings and piano''}} \\
\midrule
1 & Sustained strings melody, low quality recording     & 0.407 \\
2 & Classical wide mellow strings melody                & 0.402 \\
3 & Mellow piano melody, low quality recording          & 0.399 \\
\bottomrule
\end{tabular}
\end{table}

% ─────────────────────────────────────────────────────────────────────────────
\section{Conclusion}
% ─────────────────────────────────────────────────────────────────────────────

A two-phase fine-tuning strategy is applied to CLAP for semantic music retrieval.
Phase~1 (LoRA on RoBERTa, 0.62\% of parameters) raises MusicCaps text--audio cosine
similarity by +18.3\%. Phase~2 unfreezes the last Swin stage of the audio encoder and
trains end-to-end on raw waveforms to improve intra-genre cluster cohesion --- a metric
that Phase~1 cannot affect since it leaves the audio encoder frozen.

The image-to-playlist bridge (CLIP $\to$ CLAP MLP) is trained but not yet integrated
into the end-to-end retrieval pipeline; this is the next step.

% ─────────────────────────────────────────────────────────────────────────────
\begin{thebibliography}{9}

\bibitem{elizalde2023clap}
B.~Elizalde, S.~Deshmukh, M.~Al~Ismail, H.~Wang.
\textit{CLAP: Learning Audio Concepts from Natural Language Supervision}.
ICASSP 2023.

\bibitem{radford2021clip}
A.~Radford, J.~W.~Kim, C.~Hallacy \emph{et al.}
\textit{Learning Transferable Visual Models From Natural Language Supervision}.
ICML 2021.

\bibitem{hu2022lora}
E.~J.~Hu, Y.~Shen, P.~Wallis \emph{et al.}
\textit{LoRA: Low-Rank Adaptation of Large Language Models}.
ICLR 2022.

\bibitem{defferrard2017fma}
M.~Defferrard, K.~Benzi, P.~Vandergheynst, X.~Bresson.
\textit{FMA: A Dataset for Music Analysis}.
ISMIR 2017.

\bibitem{agostinelli2023musiclm}
A.~Agostinelli, T.~I.~Denk, Z.~Borsos \emph{et al.}
\textit{MusicLM: Generating Music From Text}.
arXiv:2301.11325, 2023.

\end{thebibliography}

\end{document}
